% ============================================
% Johns Hopkins Letter Template
% Assistant Professor Cover Letter – Muhlenberg College
% ============================================

\documentclass[11pt,letterpaper]{letter}

\usepackage{graphicx}
\usepackage{eso-pic}
\usepackage{geometry}
\usepackage{microtype}
\usepackage[hidelinks]{hyperref}

% ----- Page Setup -----
\geometry{left=25mm,right=25mm,top=25mm,bottom=25mm}

% ----- Logo Placement (first page only) -----
\newcommand{\PutJHULogoFG}[1]{%
  \AddToShipoutPictureFG*{%
    \AtPageUpperLeft{%
      \hspace{10mm}%
      \raisebox{-38mm}{%
        \includegraphics[width=6cm]{#1}%
      }%
    }%
  }%
}

% ----- Sender -----
\signature{Decory Edwards}
\address{Department of Economics \\ Johns Hopkins University \\ Baltimore, MD 21218 \\ \texttt{dedwar65@jh.edu}}

\begin{document}
\PutJHULogoFG{Images/jhulogo.png}

\begin{letter}{Search Committee \\
Accounting, Business, Economics, and Finance Department \\
Muhlenberg College \\
Allentown, PA 18104 \\ USA}

\opening{Dear Members of the Search Committee,}

I am writing to express my interest in the tenure-track Assistant Professor position in Economics/Finance in the ABEF Department at Muhlenberg College, beginning in Fall 2026. I am a Ph.D. candidate in Economics at Johns Hopkins University, advised by Professors Christopher D.~Carroll, Nicholas W.~Papageorge, and Stelios Fourakis, and I expect to receive my degree in May 2026. My research fields are macroeconomics, wealth and income dynamics, and computational economics.

My research agenda focuses on understanding the determinants of wealth inequality and the mechanisms through which heterogeneity in returns to wealth shapes aggregate outcomes. In my job-market paper, I estimate the degree of heterogeneity in individual portfolio returns using micro-data from the Survey of Consumer Finances and simulate a structural model in which households face idiosyncratic return risk. I show that allowing for persistent heterogeneity in returns substantially improves the model’s ability to match the empirical distribution of wealth, offering a tractable way to connect micro-level return dispersion to macroeconomic inequality.

In a second project, I use the Health and Retirement Study to examine the relationship between interpersonal trust and portfolio performance. Building on the literature that finds a hump-shaped relationship between trust and income, I show that a similar relationship holds for individual returns to wealth measured in the HRS data, highlighting a behavioral channel through which social attitudes shape saving and investment outcomes. This work also provides new estimates of the persistent component of returns to net worth, which motivate the calibration of heterogeneity in my structural model.

Muhlenberg’s emphasis on outstanding undergraduate teaching, innovative pedagogy, and a commitment to equity and belonging aligns strongly with my approach to both teaching and mentorship. I am particularly drawn to the College’s focus on supporting students through inclusive classrooms, interdisciplinary collaboration, and close faculty–student engagement. The ABEF Department’s expectation that research informs high-quality instruction resonates with my own view that empirical and computational tools should enrich the educational experience.

\textbf{Teaching.} My teaching experience centers on undergraduate macroeconomics and an upper-level macro policy seminar, where I have led weekly discussion sections and held regular office hours for classes ranging from 16 to 40 students. I prioritize helping students “convince themselves” that they have progressed from confusion to understanding through structured guidance and active learning. At Muhlenberg, I am prepared to teach quantitative courses across both the economics and finance majors, including \textit{Econometrics}, \textit{Principles of Microeconomics}, \textit{Principles of Macroeconomics}, and additional courses aligned with departmental needs such as corporate finance, public finance, mergers and acquisitions, or industrial organization. I look forward to developing courses that use data-intensive methods and reproducible workflows, and to involving undergraduates directly in applied research.

I believe I would be a strong fit for Muhlenberg College because my empirical and computational background allows me to (i) teach quantitative and data-driven courses effectively, (ii) integrate modern research tools such as Python, Stata, and R into classroom instruction, and (iii) maintain a research agenda that informs my teaching while involving motivated students. I am especially committed to contributing to Muhlenberg’s culture of belonging by fostering inclusive learning environments, mentoring students from all backgrounds, and supporting the College’s mission of preparing students for lives of meaning and leadership.

I would be honored to join the Muhlenberg College community and contribute to the ABEF Department's teaching, research, and service missions. Thank you for your consideration; I welcome the opportunity to discuss my fit with the department at your convenience.

\closing{Sincerely,}

\end{letter}
\end{document}
